% !TEX root = ../main.tex
\chapter{Cinématique}

% A solved problem with a verification and a figure (decimal commas in the
% axis ticks come from language=french in mytheme.sty).
\begin{problem}[title={Chute libre}]
{
Un objet est lâché sans vitesse initiale d'une hauteur $h = 20\,\mathrm{m}$.
En négligeant la résistance de l'air, combien de temps dure sa chute ?
On prend $g = 9,8\,\mathrm{m/s^2}$.
}
{
La hauteur de l'objet au temps $t$ est $y(t) = h - \frac{1}{2} g t^2$. Il
touche le sol quand $y(t) = 0$, soit
\begin{equation}
    t = \sqrt{\frac{2h}{g}} = \sqrt{\frac{2 \cdot 20}{9,8}}\,\mathrm{s}
      \approx 2,02\,\mathrm{s}.
\end{equation}
}
{
$t \approx 2,0\,\mathrm{s}$.
}
[
$\frac{1}{2} \cdot 9,8 \cdot 2,02^2 \approx 20,0\,\mathrm{m}$ : on retrouve
la hauteur de départ, et la courbe de la figure~\ref{fig:chute} coupe l'axe
du temps à $t \approx 2,02\,\mathrm{s}$.
\begin{figure}[H]
    \centering
    \includegraphics{Figures/Chute.pdf}
    \caption{Hauteur de l'objet en fonction du temps.}
    \label{fig:chute}
\end{figure}
]
\end{problem}

% Leave the solution empty: the problem is marked « À faire ».
\begin{problem}[title={Freinage}]
{
Une voiture roulant à $90\,\mathrm{km/h}$ freine avec une décélération
constante de $5,0\,\mathrm{m/s^2}$. Quelle est sa distance de freinage ?
}{}{}
\end{problem}
